\documentclass[10pt,a4paper]{amsart}
\usepackage{amsmath,amssymb,amsthm,graphicx}
\usepackage{verbatim}
\newenvironment{reference}{\comment}{\endcomment}
%\newenvironment{reference}{}{}
\newenvironment{slogan}{\comment}{\endcomment}
\newenvironment{history}{\comment}{\endcomment}

% For commutative diagrams we use Xy-pic
\usepackage[all]{xy}

% We use 2cell for 2-commutative diagrams.
\xyoption{2cell}
% Theorem environments.
%
\theoremstyle{plain}
\newtheorem{theorem}[subsection]{Theorem}
\newtheorem{proposition}[subsection]{Proposition}
\newtheorem{lemma}[subsection]{Lemma}

\theoremstyle{definition}
\newtheorem{definition}[subsection]{Definition}
\newtheorem{example}[subsection]{Example}
\newtheorem{exercise}[subsection]{Exercise}
\newtheorem{situation}[subsection]{Situation}

\theoremstyle{remark}
\newtheorem{remark}[subsection]{Remark}
\newtheorem{remarks}[subsection]{Remarks}

\numberwithin{equation}{subsection}

% Macros
%
\def\lim{\mathop{\mathrm{lim}}\nolimits}
\def\colim{\mathop{\mathrm{colim}}\nolimits}
\def\Spec{\mathop{\mathrm{Spec}}}
\def\Hom{\mathop{\mathrm{Hom}}\nolimits}
\def\Ext{\mathop{\mathrm{Ext}}\nolimits}
\def\SheafHom{\mathop{\mathcal{H}\!\mathit{om}}\nolimits}
\def\SheafExt{\mathop{\mathcal{E}\!\mathit{xt}}\nolimits}
\def\Sch{\mathit{Sch}}
\def\Mor{\mathop{\mathrm{Mor}}\nolimits}
\def\Ob{\mathop{\mathrm{Ob}}\nolimits}
\def\Sh{\mathop{\mathit{Sh}}\nolimits}
\def\NL{\mathop{N\!L}\nolimits}
\def\CH{\mathop{\mathrm{CH}}\nolimits}
\def\proetale{{pro\text{-}\acute{e}tale}}
\def\etale{{\acute{e}tale}}
\def\QCoh{\mathit{QCoh}}
\def\Ker{\mathop{\mathrm{Ker}}}
\def\Im{\mathop{\mathrm{Im}}}
\def\Coker{\mathop{\mathrm{Coker}}}
\def\Coim{\mathop{\mathrm{Coim}}}

% Boxtimes
%
\DeclareMathSymbol{\boxtimes}{\mathbin}{AMSa}{"02}

%
% Macros for moduli stacks/spaces
%
\def\QCohstack{\mathcal{QC}\!\mathit{oh}}
\def\Cohstack{\mathcal{C}\!\mathit{oh}}
\def\Spacesstack{\mathcal{S}\!\mathit{paces}}
\def\Quotfunctor{\mathrm{Quot}}
\def\Hilbfunctor{\mathrm{Hilb}}
\def\Curvesstack{\mathcal{C}\!\mathit{urves}}
\def\Polarizedstack{\mathcal{P}\!\mathit{olarized}}
\def\Complexesstack{\mathcal{C}\!\mathit{omplexes}}
% \Pic is the operator that assigns to X its picard group, usage \Pic(X)
% \Picardstack_{X/B} denotes the Picard stack of X over B
% \Picardfunctor_{X/B} denotes the Picard functor of X over B
\def\Pic{\mathop{\mathrm{Pic}}\nolimits}
\def\Picardstack{\mathcal{P}\!\mathit{ic}}
\def\Picardfunctor{\mathrm{Pic}}
\def\Deformationcategory{\mathcal{D}\!\mathit{ef}}
\usepackage{cleveref}
\crefname{originalreference}{original reference}{original references}
\Crefname{originalreference}{Original reference}{Original references}
\makeatletter
\newlabel{lemma-construct-profinite}{{097D}{source}{}{sourcepointer-1}{}}
\newlabel{lemma-construct-profinite@cref}{{[originalreference][0][]097D}{[1][0][]source}{}{}{}}
\makeatother
\crefname{originalreference}{original reference}{original references}
\Crefname{originalreference}{Original reference}{Original references}
\makeatletter
\newlabel{lemma-ind-zariski-implies}{{096T}{source}{}{sourcepointer-2}{}}
\newlabel{lemma-ind-zariski-implies@cref}{{[originalreference][0][]096T}{[1][0][]source}{}{}{}}
\makeatother
\crefname{originalreference}{original reference}{original references}
\Crefname{originalreference}{Original reference}{Original references}
\makeatletter
\newlabel{lemma-local-isomorphism-fully-faithful}{{096L}{source}{}{sourcepointer-3}{}}
\newlabel{lemma-local-isomorphism-fully-faithful@cref}{{[originalreference][0][]096L}{[1][0][]source}{}{}{}}
\makeatother
\crefname{originalreference}{original reference}{original references}
\Crefname{originalreference}{Original reference}{Original references}
\makeatletter
\newlabel{lemma-silly}{{00DS}{source}{}{sourcepointer-4}{}}
\newlabel{lemma-silly@cref}{{[originalreference][0][]00DS}{[1][0][]source}{}{}{}}
\makeatother
\crefname{originalreference}{original reference}{original references}
\Crefname{originalreference}{Original reference}{Original references}
\makeatletter
\newlabel{lemma-w-local-morphism-equal-points-stalks-is-iso}{{097E}{source}{}{sourcepointer-5}{}}
\newlabel{lemma-w-local-morphism-equal-points-stalks-is-iso@cref}{{[originalreference][0][]097E}{[1][0][]source}{}{}{}}
\makeatother
\title{A w-local morphism of w-local affine schemes which identifies local rings is ind-Zariski}
\author{The Stacks Project Authors}
\date{stacks-project:proetale.tex:local-snapshot-20261012; GNU FDL 1.2 or later}
\begin{document}
\maketitle
\noindent\textit{Editorial note.} All delivered mathematical text is the author's own, copied from the pinned Stacks Project chapter file proetale.tex; no statement or proof step was rewritten, reordered or removed.\par\medskip
\noindent\textit{Editorial note.} The counted proof is the author's complete argument: the construction of the ind-Zariski cover with prescribed components, the passage from the continuous map on components to a morphism of affine schemes by faithfulness of local isomorphisms, and the closing isomorphism criterion for w-local morphisms.\par\medskip
\noindent\textit{Editorial note.} The label lemma-w-local-morphism-equal-stalks-is-ind-zariski names the counted statement, and the proof environment immediately following it in the source is delivered verbatim as its proof.\par\medskip
\noindent\textit{Editorial note.} The counted statement is carried by the author's own lemma environment in the chapter source.\par\medskip
\noindent\textit{Editorial note.} Pre-open gate: python operations/agent-authoring/tools/precard\_gate.py --source proetale.tex --label lemma-w-local-morphism-equal-stalks-is-ind-zariski --start 1184 --end 1221 --local-id 3594 returned RESULT=OK (L3594 free, occupancy FREE, public interval clear). Candidate taken from sources/stacks-proetale/candidates-halfA.json.\par\medskip
\noindent\textit{Editorial note.} The statement and proof are copied from the pinned chapter file; the cited inputs are kept as external pointers with their canonical Stacks tags.\par\medskip
\noindent\textit{Allowed standard prerequisites (editorial).}\begin{itemize}
\item Realization of a profinite space of components by an ind-Zariski cover. Same chapter, the lemma producing an ind-Zariski ring map realizing a prescribed profinite space of connected components. Supplies the ind-Zariski map A -> A' with the same connected components as B and with cartesian component diagram.
\item Identification of the closed points of the constructed cover. Same chapter, the lemma describing the closed points of the cover as the inverse image of the closed points of the base. Supplies the w-locality of X' and the description of its closed points.
\item Faithfulness of the space map for local isomorphisms, and the isomorphism criterion. Same chapter, the lemmas that a local isomorphism is fully faithful and that a w-local morphism of w-local schemes which is bijective on closed points is an isomorphism. Supply the morphism of affine schemes Y -> X' over X corresponding to the continuous map, and the conclusion that it is an isomorphism.
\end{itemize}\medskip
\ifdefined\hypertarget\hypertarget{sourcepointer-1}{}\fi
\noindent\textit{Source reference.} Original reference 097D: sources/stacks-proetale/pinned-original-source.tex, source lines 1098--1098. Prerequisite pointer invoked at proof line(s) 1207; it names a standard fact of the same chapter that the counted argument uses, and every step of the counted proof built on it is retained verbatim in the delivered body.\par\medskip
\ifdefined\hypertarget\hypertarget{sourcepointer-2}{}\fi
\noindent\textit{Source reference.} Original reference 096T: sources/stacks-proetale/pinned-original-source.tex, source lines 685--685. Prerequisite pointer invoked at proof line(s) 1215; it names a standard fact of the same chapter that the counted argument uses, and every step of the counted proof built on it is retained verbatim in the delivered body.\par\medskip
\ifdefined\hypertarget\hypertarget{sourcepointer-3}{}\fi
\noindent\textit{Source reference.} Original reference 096L: sources/stacks-proetale/pinned-original-source.tex, source lines 603--603. Prerequisite pointer invoked at proof line(s) 1214; it names a standard fact of the same chapter that the counted argument uses, and every step of the counted proof built on it is retained verbatim in the delivered body.\par\medskip
\ifdefined\hypertarget\hypertarget{sourcepointer-4}{}\fi
\noindent\textit{Source reference.} Original reference 00DS: sources/stacks-proetale/pinned-original-source.tex, source lines 434--434. Prerequisite pointer invoked at proof line(s) 1208; it names a standard fact of the same chapter that the counted argument uses, and every step of the counted proof built on it is retained verbatim in the delivered body.\par\medskip
\ifdefined\hypertarget\hypertarget{sourcepointer-5}{}\fi
\noindent\textit{Source reference.} Original reference 097E: sources/stacks-proetale/pinned-original-source.tex, source lines 1155--1155. Prerequisite pointer invoked at proof line(s) 1219; it names a standard fact of the same chapter that the counted argument uses, and every step of the counted proof built on it is retained verbatim in the delivered body.\par\medskip
\begin{lemma}
\label{lemma-w-local-morphism-equal-stalks-is-ind-zariski}
Let $A \to B$ be ring map such that
\begin{enumerate}
\item $A \to B$ identifies local rings,
\item the topological spaces $\Spec(B)$, $\Spec(A)$ are w-local, and
\item $\Spec(B) \to \Spec(A)$ is w-local.
\end{enumerate}
Then $A \to B$ is ind-Zariski.
\end{lemma}

\begin{proof}
Set $X = \Spec(A)$ and $Y = \Spec(B)$. Let $X_0 \subset X$ and
$Y_0 \subset Y$ be the set of closed points. Let $A \to A'$ be the ind-Zariski
morphism of affine schemes such that with $X' = \Spec(A')$ the diagram
$$
\xymatrix{
X' \ar[r] \ar[d] & \pi_0(X') \ar[d] \\
X \ar[r] & \pi_0(X)
}
$$
is cartesian in the category of topological spaces and such that
$\pi_0(X') = \pi_0(Y)$ as spaces over $\pi_0(X)$, see
Lemma \ref{lemma-construct-profinite}. By
Lemma \ref{lemma-silly} we see that $X'$ is w-local and
the set of closed points $X'_0 \subset X'$ is the inverse image of $X_0$.

\medskip\noindent
We obtain a continuous map $Y \to X'$ of underlying topological spaces
over $X$ identifying $\pi_0(Y)$ with $\pi_0(X')$. By
Lemma \ref{lemma-local-isomorphism-fully-faithful}
(and Lemma \ref{lemma-ind-zariski-implies})
this corresponds to a morphism of affine schemes $Y \to X'$
over $X$. Since $Y \to X$ maps $Y_0$ into $X_0$ we see that
$Y \to X'$ maps $Y_0$ into $X'_0$, i.e., $Y \to X'$ is w-local.
By Lemma \ref{lemma-w-local-morphism-equal-points-stalks-is-iso}
we see that $Y \cong X'$ and we win.
\end{proof}

\end{document}
